\section{Relasi Logika Bersyarat: Implikasi dan Biimplikasi}

Setelah mencicipi ketegasan alat pelekat dasar, ranah logika kian meruncing menyekat wilayah para \textit{Programmer} jagoan melalui pisau pamungkas hierarki tatanan argumen: Pernyataan bersyarat pemicu rantai kejadian masa depan. Tak lain duo kembar ningrat \textbf{Implikasi Bersyarat} dan perwujudan final cermin ganda ekuivalennya: \textbf{Bikondisional / Biimplikasi} \cite{rosen2019}.

\subsection{1. Implikasi: Janji Mengancam "Jika... Maka..." ($\rightarrow$)}

Struktur tulang punggung dari percabangan perintah \textit{Syntax koding "IF ... THEN ..."} di bahasa C, Python, maupun Java lahir berawal dari rumus absolut biner Implikasi $p \rightarrow q$.

Si pencetus pemicu bom panah kiri $p$ mulia dititah gelar keraton "Hipotesis / Anteseden", sedangkan wujud anak cucu akibat ledakan panah penerima kanan $q$ menyandang takhta "Konklusi / Konsekuen".

\begin{definisi}[Perangai Aneh Implikasi]
Tabel penyusunan takdir nilai sang Implikasi acapkali melahirkan protes kejanggalan nalar kaum natural awam! Paksaan pilar $p \rightarrow q$ \textbf{HANYA DIKUTUK MENETAP PALU SALAH (F) PADA SATU TITIK NASIB TUNGGAL MANAKALA:} Pemicu asalnya dengan ceroboh telah matang terjadi utuh merona Benar (T), Namun si akibat yang dijanjikannya ingkar janji cacat parah terbukti memble di nilai False (F). Selain 1 noda pelanggaran cacat itu, barisan sisa 3 kombinasi semu lainnya dianugerahi kemuliaan pahala \textbf{BENAR MUTLAK (T)} walau si pemicunya tak terjadi sama sekali di awal dongeng!
\end{definisi}

\textbf{Contoh Anomali Tabel yang Membingungkan (Klausul \textit{Vacuous Truth}):} \\
Narasi Caleg Pilkada korup: *"Bila saya fiks terpilih memenangkan tiket takhta kursi Gubernur esok lusa (p), saya mengikat Janji Darah bakal menyebar membagikan miliaran laptop Macbook Pro kepada mahasiswa sekota raya secara gratis (q)!"*

Evaluasi Hukum Pemilu: 
\begin{itemize}
    \item \textbf{Barisan Paling Normal (T $\rightarrow$ T):} Sang Bapak korup tembus memenangkan jengkal kursi kekuasaan gubernur. Beliau kebetulan lagi dermawan pamer waras mengirim kargo raksasa menyebarkan jutaan keping Macbook ke asrama kos mahasiswa persis menepati janji. Apakah janji hukum \textit{Implikasi} dinobatkan berhasil dan sah nilainya? **YA MUTLAK BENAR (T)!** (Status aman sudi tepati komitmen).
    \item \textbf{Barisan Jatuh Cacatan Pelanggaran Maut Tunggal (T $\rightarrow$ F):} Terkutuklah Sang Bapak malang terpilih sukses resmi duduk di kursi marmer gubernur, Eh namun tak membagikan setetes laptop sebiji kawatpun di malam kampusnya. Sang penagih janji (Mahasiswa) menuntut demo merobohkan pagar menitis kejam. Janjinya dituduh **BOHONG INGKARI SUMPAH PALSU CACAT SALAH MUTLAK (F)!**
    \item \textbf{Barisan Ganjil "Kebenaran Vakum Hampa" (F $\rightarrow$ entah T/F):} Realitas menangis kejam sang jagoan bapak caleg ini cuma kantongi receh suara memble Kalah babak belur di pilkada Gubernur! Di bulan kedelapan lantaran asyik terlewat galau sedihnya ia iseng malah menenggelamkan diri menguras atm melempar beli Macbook di-donasikan iseng mendadak ke kampus antah berantah. Pertanyaan ajaib hakim Hukum: "Lantas, apakah janji ikatan orisinal Implikasinya sejak masa awal bulan pemilu DULU ITU tercoreng dinilai pembohong cacat secara hukum aljabar persidangan?!" Jawabannya \textbf{TIDAK! Janjinya dinilai Absolut Murni tetap sah menyandang panji BENAR (T)} di atas dokumen. Sang bapak bukan pelanggar pasal penipuan janji. *Wong syarat letupan komando pemicu awal gubernurnya gagal menduduki tahta \textbf{(F)}, alam semesta memaafkannya bebas tak ditagih pertanggungjawaban panah q akhir mau diledakkan macem pamer apapun juga di ujung waktu itu terserah \textbf{(Prinsip Sakral Vacuous Truth)}.*
\end{itemize}

\subsection{2. Biimplikasi: Bayangan Cermin Ganda Setara ($\leftrightarrow$)}

Manakala dua kelompok geng mafia geng motor memaksakan sebuah perjanjian \textit{Treaty} syarat ganda berat ketimbang sekedar klausul sepihak implikasi kasta atas. Sang \textit{Bikondisional / Biimplikasi} lahir mengunci gerbang relasi 2 arah di simbol $p \leftrightarrow q$ yang ditakbirkan *"p Berlaku bilamana, dan hanya sejatinya bilamana, diiringi kepatuhan q."*

\begin{definisi}[Sifat Absolut Biimplikasi Cermin Ganda]
Kompak bersiul di relasi ekuivalen murni tak tertembus asimetris! Proposisi Majemuk Biimplikasi bertakhta mahkota **T (BENAR)** HANYA di kala status pilar pemicu kanan kembar dempet identik cermin merangkul seiya sekata merunut absolut memeluk kodrat kasta milik pilar panah subjek di kubu sisi kompas kiri asalnya (Sama-sama Serentak T. Ataukah malah kompak menyeburkan diri Setia kawan merana berdua False). Setengah belang melanggar sedikit setetes serpihan sumbang (T \& F diacak) mutlak merobohkan istana menelan kebohongan FALSE raksasa tak berampun.
\end{definisi}

Contoh Kasus Eksekusi Biimplikasi Syarat Mutlak Ganda di dunia Server IT: \\
"Identitas \textit{Token} \textit{Password Mobile M-Banking} Anda \textbf{SUCCES} dipindai mesin server otentikasi lolos divalidasi pintu \textbf{JIKA DAN BILA DAN HANYA JIKA BILA} sistem rekam biometrik layar handphone ini sungguh memindai mendeteksi 100\% retensi pola sidik jari jempol sang tuan aslinya mencocok persis rapot." \\
(Gagal di sidik jari jempol F padahal password berhasil dilolos T? Biimplikasi Bank meretas palu *Fraud Alert F* memblokir saldo ratusan juta seketika tanpa kompromi!) \cite{wiki_first_order_logic}.
